\chapter{CMake 与依赖}
\index{CMake Presets}\index{dependency management}\index{compile commands@compile\_commands.json}

\begin{introduction}
  \item 区分配置、生成、构建与测试，读懂源目录和构建目录
  \item 用库、应用、示例和测试组成 target-based CMake 构建图
  \item 准确选择 \texttt{PRIVATE}、\texttt{PUBLIC} 与 \texttt{INTERFACE}
  \item 从干净目录复现 Debug、Sanitized 与 Release 配置
\end{introduction}

\chaptermeta{已完成第 1、3、5 章，能编译多文件程序并理解头文件与链接。}{Python 的 \path{pyproject.toml} 主要描述包和工具；CMake 描述原生目标、编译要求与目标之间的依赖边。}{先确认失败发生在配置、生成、编译、链接、测试哪一步，再检查目标图；不要用全局 include 路径掩盖缺失依赖。}

\section{本章输出：从源码树复现一个目标图}

量化程序从一个可执行文件增长为公共计算库、命令行应用、研究示例和行为测试后，继续复制编译选项与源文件列表会产生漂移。本章把一个 VWAP 汇总器组织成六个正常目标：\texttt{ch11\_project\_options} 保存共同构建要求，\texttt{ch11\_market} 提供静态库，报告程序、示例、测试和练习答案只声明自己对该库的依赖。另有一个 \texttt{EXCLUDE\_FROM\_ALL} 故障目标，仅供显式触发链接实验，不属于正常构建。

从项目根目录执行以下命令；\path{build/ch11-debug} 必须在开始前不存在：

\begin{lstlisting}[language=bash,numbers=none]
cmake -S code/ch11 -B build/ch11-debug -DCMAKE_BUILD_TYPE=Debug
cmake --build build/ch11-debug -j
ctest --test-dir build/ch11-debug --output-on-failure
./build/ch11-debug/ch11_market_report
\end{lstlisting}

最后一条命令精确输出：

\begin{lstlisting}[numbers=none]
rows=3 quantity=25 notional=2525.00 vwap=101.00
\end{lstlisting}

完成判据不只是“我的机器能编译”：全新构建树能配置并构建全部目标；CTest 运行库的公开行为；应用无需自己声明头文件目录或 C++ 标准；移除链接边会稳定停在链接阶段；Debug、Sanitized 和 Release 使用彼此独立的构建树。

\section{配置不是编译：先读四个阶段}

\subsection{场景：为什么旧可执行文件会制造假绿}

\texttt{cmake -S ... -B ...} 先读取 \path{CMakeLists.txt}、检测编译器和依赖，再生成 Ninja 或 Make 等底层构建文件。它通常不编译 C++。\texttt{cmake --build ...} 才按依赖图调用编译器与链接器；\texttt{ctest --test-dir ...} 运行已经登记且已经构建的测试。若只运行最后一步，旧二进制可能让错误源码看起来仍然通过。

CMakeCache.txt 保存编译器路径、选项和检测结果，所以仅删除一个目标文件无法得到干净配置。验证可复现性时应换用从未存在的构建目录；切换编译器、生成器或 sanitizer 时也优先使用新目录，反复修改同一个缓存容易残留旧状态。

\begin{pythonbridge}
Python 直接执行源码时，安装与运行边界常被解释器隐藏；原生项目把配置、生成、编译、链接和测试明确分开。CMake 更接近“生成构建系统的项目模型”，既不承担编译器职责，也不能等同于 pip 包管理器。先判断命令停在哪个阶段，才能决定检查 CMake 目标、C++ 源码、链接库还是运行行为。

\pythoncompare{Python 常直接执行源码；CMake 显式分开配置、构建、链接和测试。}{Python 环境拥有包；原生 target 拥有编译选项和传递依赖。}{Python 导入时才暴露部分问题；CMake 生成和编译阶段分层检查。}{失败分别落在配置、编译、链接或运行，不能统称为构建失败。}
\end{pythonbridge}

\section{用目标表达库、应用、测试与示例}

先观察重构顺序。第一步若只有报告程序，常见写法是把应用和实现源文件都塞给一个 executable：

\begin{lstlisting}[language=bash,numbers=none]
add_executable(ch11_market_report
  app/main.cpp src/market_summary.cpp)
target_include_directories(ch11_market_report PRIVATE include)
\end{lstlisting}

它能构建，却没有可供示例和测试复用的库。第二步抽出 \texttt{ch11\_market}，把实现源文件和公共 include 要求移到库目标；报告程序删除重复源文件，只保留一条链接边：

\begin{lstlisting}[language=bash,numbers=none]
add_library(ch11_market STATIC src/market_summary.cpp)
target_include_directories(ch11_market PUBLIC include)
add_executable(ch11_market_report app/main.cpp)
target_link_libraries(ch11_market_report PRIVATE ch11_market)
\end{lstlisting}

第三步按相同形状增加 example 与 test：各自拥有入口源码，各自 PRIVATE 链接库；测试再用 \texttt{add\_test} 登记。最后才把多目标共同的 C++20、警告和 sanitizer 要求收进 INTERFACE options 目标。每一步都删除旧目标上的重复属性，并立即从空构建目录验证，避免最终清单看似整齐却无法解释依赖从何而来。

下面的完整清单就是独立章节快照的构建入口：

\lstinputlisting[caption={第 11 章完整 target-based CMake 项目},language=bash]{code/ch11/CMakeLists.txt}

\texttt{add\_library} 与 \texttt{add\_executable} 创建有名称的目标。源文件、头文件目录、语言标准、编译定义和链接项随后附着到具体目标，不写成影响整个目录的全局状态。应用、示例、测试都只通过 \texttt{target\_link\_libraries(... ch11\_market)} 建立一条边；它们不复制 \path{src/market_summary.cpp}，因此库的实现只有一个编译实体。

\texttt{include(CTest)} 提供 \texttt{BUILD\_TESTING} 约定并启用测试登记；\texttt{add\_test} 描述运行期行为，但不会自动创建或构建可执行文件。测试本身仍是普通目标，区别只在于它被 CTest 调度。

公共头文件展示了库消费者真正需要的 API：

\lstinputlisting[caption={由库目标公开的行情汇总接口}]{code/ch11/include/quant/ch11/market_summary.hpp}

该接口使用 \texttt{std::span}，所以消费者也必须以 C++20 编译。应让库目标传播这一使用要求，免去每个应用分别记住 \texttt{-std=c++20}。

\section{PRIVATE、PUBLIC 与 INTERFACE 是传播方向}

三个关键字必须针对“某一项目标属性”理解，而不是给整个库贴标签：

\begin{center}
\begin{tabular}{lll}
\toprule
关键字 & 当前目标使用 & 链接当前目标的消费者使用 \\
\midrule
\texttt{PRIVATE} & 是 & 否 \\
\texttt{PUBLIC} & 是 & 是 \\
\texttt{INTERFACE} & 否 & 是 \\
\bottomrule
\end{tabular}
\end{center}

\path{include/} 中的头文件是公共接口，所以 \texttt{target\_include\_directories(ch11\_market PUBLIC include)} 既让库实现找到头文件，也让应用找到它。\texttt{ch11\_project\_options} 没有源码、不会生成库文件；它是 INTERFACE 目标，只携带 C++20、警告和可选 sanitizer 要求。市场库以 PUBLIC 链接它，应用因而沿目标边获得这些要求。

相反，\texttt{CH11\_BUILDING\_LIBRARY} 只用于编译库实现，故声明为 PRIVATE。应用源码包含一个反向检查：

\lstinputlisting[caption={消费者使用公共接口且拒绝私有定义泄漏}]{code/ch11/app/main.cpp}

若把该编译定义误改为 PUBLIC，应用中的 \texttt{\#error} 会在编译阶段触发。这是传播行为的直接证据，不需要读取 CMake 内部属性。真实可安装库还会用 \texttt{BUILD\_INTERFACE} 和 \texttt{INSTALL\_INTERFACE} 区分源码树与安装后的头文件位置；本章快照只在源码树构建，暂不引入安装规则。

失败诊断：将报告目标的 \texttt{target\_link\_libraries} 行删除。头文件可能先因 include 路径不再传播而编译失败；若手工加全局 include 路径，编译可能通过，最终仍因 \texttt{summarize} 没有链接定义而失败。这说明“加一个 \texttt{-I}”不能修复缺失的目标边。

仓库还隔离了一个最小链接实验：

\lstinputlisting[caption={声明存在但没有目标提供定义的故障样例}]{code/ch11/failures/missing_target_link.cpp}

\begin{lstlisting}[language=bash,numbers=none]
g++ -std=c++20 code/ch11/failures/missing_target_link.cpp \
  -o build/missing-target-link
\end{lstlisting}

源文件能通过编译，但链接器报告 undefined reference；恢复方法是链接提供定义的目标，而不是在声明后再写一份临时实现。

\section{Debug、Sanitized 与 Release 各自回答什么}

单配置生成器通过 \texttt{CMAKE\_BUILD\_TYPE} 选择一组配置。Debug 保留调试信息并适合逐步定位；Sanitized 在 Debug 基础上增加 AddressSanitizer 与 UndefinedBehaviorSanitizer，用运行期开销换取内存和 UB 证据；Release 启用优化并常定义 \texttt{NDEBUG}，适合最终行为与性能验收，不能代替前两者。

\begin{lstlisting}[language=bash,numbers=none]
cmake -S code/ch11 -B build/ch11-debug -DCMAKE_BUILD_TYPE=Debug
cmake --build build/ch11-debug -j
ctest --test-dir build/ch11-debug --output-on-failure

cmake -S code/ch11 -B build/ch11-sanitized \
  -DCMAKE_BUILD_TYPE=Debug -DCH11_ENABLE_SANITIZERS=ON
cmake --build build/ch11-sanitized -j
ctest --test-dir build/ch11-sanitized --output-on-failure

cmake -S code/ch11 -B build/ch11-release -DCMAKE_BUILD_TYPE=Release
cmake --build build/ch11-release -j
ctest --test-dir build/ch11-release --output-on-failure
\end{lstlisting}

三个目录并存，使配置差异可审计，也避免缓存串味。Sanitizer 编译选项与链接选项都挂在 INTERFACE 目标上：只加编译选项可能在最终链接 sanitizer runtime 时失败。MSVC 分支使用 \texttt{/W4 /permissive-}；本章的 sanitizer 开关只对 GCC/Clang 主线生效。

多配置生成器（如 Visual Studio）通常在一次配置中生成 Debug 与 Release，\texttt{CMAKE\_BUILD\_TYPE} 会被忽略；构建和测试时写 \texttt{cmake --build build --config Release} 与 \texttt{ctest --test-dir build -C Release}。MinGW Makefiles 通常仍是单配置生成器，沿用 \texttt{CMAKE\_BUILD\_TYPE} 和 GCC 风格选项，可执行文件带 \path{.exe}；它不能与 MSVC 编译出的不兼容 ABI 库随意混用。平台改变了选择配置的命令位置和产物后缀，没有改变目标、属性和依赖边的模型。

\section{Presets、clangd 与 CI 共享同一配置事实}

根目录的 \path{CMakePresets.json} 把常用 Linux/WSL Debug、Linux/WSL Release 和 Windows MinGW Release 配置登记为版本化名称。它不保存开发者本机的绝对路径或密钥；这些个人覆盖应放进不提交的 \path{CMakeUserPresets.json}。在 WSL 中可直接运行：

\begin{lstlisting}[language=bash,numbers=none]
cmake --preset linux-debug
cmake --build --preset linux-debug
ctest --preset linux-debug
\end{lstlisting}

所有 preset 都打开 \texttt{CMAKE\_EXPORT\_COMPILE\_COMMANDS}，生成的 \path{compile_commands.json} 让 clangd、clang-tidy 和编辑器取得每个翻译单元真实的 include 路径、宏与语言标准。项目同时提交 \path{.clang-format} 和 \path{.clang-tidy}，但格式化与静态分析仍应在单独改动中运行并审阅；自动改写整仓不能替代编译和行为测试。

\texttt{include(CTest)} 提供 \texttt{BUILD\_TESTING} 选项。测试专用 executable 与 \texttt{add\_test} 应放在 \texttt{if(BUILD\_TESTING)} 内，使只消费库或应用的构建不必编译测试框架；正常应用、示例与 benchmark 是否保留，则由部署范围决定。\path{code/ch11/CMakeLists.txt} 和最终项目都采用这条边界。

\path{.github/workflows/ci.yml} 在 GCC/Clang 与 Debug/Release 的四个组合中执行配置、构建和 CTest，另用独立 job 验收书稿契约。矩阵的价值在于让编译器和配置成为显式变量，不在于追求格子数量；sanitizer、fuzzing 或昂贵 benchmark 应按运行时间拆成独立 job，并保留失败产物。

\section{依赖锁定、测试框架与 fuzzing 的采用边界}

Conan 或 vcpkg 都可以把外部库解析为 CMake imported target，但可复现性还需要提交 manifest、版本约束和 lockfile/baseline。二者不应同时成为同一项目的默认依赖来源。本仓库的 \path{examples/dependency-locking/vcpkg.json} 锁定 baseline，并声明 Catch2、Eigen3 和 pybind11；\path{capstones/factor_kernel} 的可选 target 则真实构建 Catch2 参数化测试与 Eigen kernel。前半部仍使用无依赖程序展示最小测试协议；后半部引入框架时，公开行为边界与独立期望值仍保持不变。

解析器适合用 coverage-guided fuzzing 补充手写样例。根项目的可选目标 \texttt{quant\_csv\_fuzzer} 只在显式设置 \texttt{-DQUANT\_BUILD\_FUZZERS=ON} 且使用 Clang/libFuzzer 时生成：

\begin{lstlisting}[language=bash,numbers=none]
CC=clang CXX=clang++ cmake -S . -B build/fuzz \
  -DQUANT_BUILD_FUZZERS=ON -DBUILD_TESTING=OFF
cmake --build build/fuzz --target quant_csv_fuzzer -j
./build/fuzz/quant_csv_fuzzer -max_total_time=60 corpus/csv
\end{lstlisting}

fuzzer 的 oracle 不是“永不报错”：任意字节输入都可以被拒绝，但成功结果必须只包含非空代码、有限正价格与正数量，而且 sanitizer 不得报告内存或未定义行为。发现崩溃后保存最小语料，并转成普通回归测试；持续随机运行不能替代稳定复现。

\section{外部依赖也应表现为目标}

现代包提供的 CMake 配置通常暴露命名空间目标。以下只是语法示意，本章不下载 fmt：

\begin{lstlisting}[language=bash,numbers=none]
find_package(fmt 10 CONFIG REQUIRED)
target_link_libraries(ch11_market_report PRIVATE fmt::fmt)
\end{lstlisting}

\texttt{fmt::fmt} 应携带自己的头文件目录、编译定义及传递链接项；消费者无需猜测安装路径。固定可接受的版本范围、记录工具链和依赖来源，才能使“我的缓存里能找到”变成可复现配置。若使用 \texttt{FetchContent}，还需固定提交或发布版本，并明确网络与离线策略；不要默认追踪上游主分支。

故障诊断按顺序进行：配置期找不到包，检查工具链文件、前缀路径和版本；编译期找不到头文件，检查是否链接了正确 imported target 及其可见性；链接期缺符号，检查库目标和 ABI/配置是否匹配；运行期找不到动态库，再检查部署路径。把所有路径塞进全局 \texttt{include\_directories} 或 \texttt{link\_directories} 会让无关目标偶然成功，并隐藏真实依赖。

\section{章末三层练习}

\subsection*{随堂验证汇总}

\begin{enumerate}
  \item 只运行配置命令，确认构建目录中已有缓存和生成文件但还没有 \texttt{ch11\_market\_report}；再构建该单一目标，观察静态库先于应用生成；最后运行 CTest，解释为何未构建测试目标时测试可能显示“找不到可执行文件”。
  \item 把 \texttt{ch11\_market} 对 \texttt{ch11\_project\_options} 的 PUBLIC 暂改为 PRIVATE，在全新构建树中观察消费者能否继续使用 \texttt{std::span}；恢复后再把私有编译定义改为 PUBLIC，确认失败位置和第一条稳定诊断不同。
  \item 在纸上为“CSV 库、回测库、命令行应用、单元测试”画目标图。逐条标出公共头文件、仅实现使用的编译定义、消费者必须继承的 C++ 标准和外部库，随后为每条属性选择 PRIVATE、PUBLIC 或 INTERFACE；若某属性没有明确消费者，先不要传播。
\end{enumerate}

复现层从干净目录构建既有 target；修改层新增下述名义金额目标；开放层为一个真实依赖写出固定版本、可替换边界与离线失败政策。

本章遵循 ADR 0001，把 \path{code/ch11} 保留为可独立配置的阶段快照。整书根项目通过 \texttt{add\_subdirectory} 复用同一目标声明；自动验收还会删除旧测试构建树，从空目录分别构建普通与 Sanitized 配置下的全部正常目标、运行该快照自己的 CTest，再核对报告程序输出。这样 PDF 清单、读者命令与 CI 使用同一份 CMake 事实源。

编码任务：新增一个 \texttt{ch11\_notional\_solution} 可执行目标。源码只包含公共头文件并链接 \texttt{ch11\_market}，不得复制库源文件或手写 include 路径。用两行行情 \texttt{\{99.0,1\}} 与 \texttt{\{101.0,1\}} 精确输出：

\begin{lstlisting}[numbers=none]
exercise-notional=200.00
\end{lstlisting}

从全新 Debug 构建树构建该单一目标并运行，再构建全项目和运行 CTest。完整源码见附录；目标声明应只需要 \texttt{add\_executable} 与一条 PRIVATE 链接边。

自测：

\begin{enumerate}
  \item 配置、生成、构建和测试分别读取什么、产生什么，哪一步调用编译器？
  \item 为什么存在旧可执行文件不能证明当前源码可从零构建？
  \item 库的公共头文件目录为何是 PUBLIC，而库实现宏为何是 PRIVATE？
  \item INTERFACE 目标既不编译源码，为何仍能影响应用的编译和链接？
  \item \texttt{add\_test} 为什么不能代替 \texttt{add\_executable}？
  \item 缺 include 路径与缺函数定义分别最可能在哪个阶段暴露？
  \item Debug、Sanitized 和 Release 各自提供什么证据，为什么使用独立构建树？
  \item 单配置与多配置生成器选择 Release 的命令差异是什么，哪些目标模型不变？
\end{enumerate}


\section{本章小结}

CMake 的核心产物是目标图。目标拥有源码和构建属性，依赖边传播真实使用要求；PRIVATE、PUBLIC、INTERFACE 描述传播方向。配置、构建和测试是不同阶段，旧缓存不能替代干净复现。Debug、Sanitized 与 Release 回答不同问题，外部依赖也应以版本化目标进入图。只要应用、测试和示例都通过链接库获得所需信息，项目结构就能随量化代码增长而保持可解释、可验证。
